\begin{proposition}[4.8.5]
\label{IV.2.4.8.5-fr}
Con las notaciones de \textup{(4.8.2)}, sea $(X_\alpha)$ un recubrimiento abierto de
$X$. Para que un elemento $\mathscr H\in\Phi(\mathscr F\otimes_kK)$ \textup{(resp.}
$u\in\operatorname{Hom}(\mathscr F\otimes_kK,\mathscr G\otimes_kK)$, \textup{resp.}
$Z\in\mathfrak S(X_{(K)})$, \textup{resp.} $M\in\mathfrak P(X_{(K)})$\textup{)} esté definido
sobre $K'$, es necesario y suficiente que, para todo $\alpha$,
$\mathscr H|(X_\alpha)_{(K)}$ \textup{(resp.} $u|(X_\alpha)_{(K)}$,
$Z\cap(X_\alpha)_{(K)}$, $M\cap(X_\alpha)_{(K)}$\textup{)} esté definido sobre $K'$.
\end{proposition}

\begin{proof}
Esto resulta inmediatamente de la inyectividad de las aplicaciones (4.8.2.1),
(4.8.2.2), (4.8.2.3) y (4.8.2.5): por ejemplo, si para todo $\alpha$ existe un
sub-$\mathcal O_{(X_\alpha)_{(K')}}$-módulo cuasicoherente
$\mathscr H'_\alpha$ de $(\mathscr F|X_\alpha)\otimes_kK'$ tal que
$\mathscr H|(X_\alpha)_{(K)}=\mathscr H'_\alpha\otimes_{K'}K$, para dos
índices cualesquiera $\alpha$, $\beta$, se tendrá
$\mathscr H|(X_\alpha\cap X_\beta)_{(K)}
=\bigl(\mathscr H'_\alpha|(X_\alpha\cap X_\beta)_{(K')}\bigr)
 \otimes_{K'}K
=\bigl(\mathscr H'_\beta|(X_\alpha\cap X_\beta)_{(K')}\bigr)
 \otimes_{K'}K$,
por tanto, necesariamente,
$\mathscr H'_\alpha|(X_\alpha\cap X_\beta)_{(K')}
=\mathscr H'_\beta|(X_\alpha\cap X_\beta)_{(K')}$ cualesquiera que sean
$\alpha$ y $\beta$, y existe, por consiguiente, un
sub-$\mathcal O_{X_{(K')}}$-módulo cuasicoherente $\mathscr H'$ de
$\mathscr F\otimes_kK'$ tal que
$\mathscr H'|(X_\alpha)_{(K')}=\mathscr H'_\alpha$ para todo $\alpha$, y por
tanto $\mathscr H=\mathscr H'\otimes_{K'}K$. Se razona del mismo modo en los
demás casos.
\end{proof}

\begin{lemma}[4.8.6]
\label{IV.2.4.8.6-fr}
Sean $k$ un cuerpo, $K$ una extensión de $k$, $A$ una $k$-álgebra, $M$ un
$A$-módulo, $\overline N$ un sub-$A_{(K)}$-módulo de $M_{(K)}$, y $K'$ una
subextensión de $K$. Las condiciones siguientes son equivalentes:
\begin{enumerate}
\item[\textup{a)}] $\overline N$ es de la forma $N'\otimes_{K'}K$, donde $N'$ es un
sub-$A_{(K')}$-módulo de $M_{(K')}$.
\item[\textup{b)}] Si $X=\operatorname{Spec}(A)$, el
$\mathcal O_{X_{(K)}}$-módulo cuasicoherente $(\overline N)^{\sim}$ sobre
$X_{(K)}=\operatorname{Spec}(A_{(K)})$ está definido sobre $K'$.
\item[\textup{c)}] $\overline N$ es de la forma $N'\otimes_{K'}K$, donde $N'$ es un
subespacio vectorial sobre $K'$ de $M_{(K')}$.
\end{enumerate}
\end{lemma}

\oldpage[IV-2]{82}

\begin{proof}
La equivalencia de \emph{a)} y \emph{b)} resulta trivialmente de (\textbf I, 1.6.5); por otra parte,
es claro que \emph{a)} implica \emph{c)}. Recíprocamente, si se cumple \emph{c)}, se sabe que puede
escribirse (salvo identificación canónica)
$N'=\overline N\cap M_{(K')}$. Como $\overline N$ y $M_{(K')}$ son, por
hipótesis, $A_{(K')}$-módulos, lo mismo ocurre con $N'$.
\end{proof}

\begin{corollary}[4.8.7]
\label{IV.2.4.8.7-fr}
Bajo las hipótesis de \textup{(4.8.6)}, existe un menor subcuerpo $K'$ de $K$ que
contiene $k$ tal que se cumplen las condiciones equivalentes de \textup{(4.8.6)}.
\end{corollary}

\begin{proof}
Basta verlo para la condición \emph{c)}, en cuyo caso resulta de Bourbaki,
\emph{Alg.}, cap.~II, 3.\textsuperscript{a} ed., \S~8,
n.\textsuperscript{o}~6, prop.~6.
\end{proof}
